\documentclass[11pt]{article}
\parindent=0in
\textwidth=6.7in 
\textheight=9.5in
\setlength{\oddsidemargin}{-0.10in}
\setlength{\evensidemargin}{-0.10in}
\setlength{\topmargin}{-0.8in}
\setlength{\unitlength}{0.5in}

\usepackage[utf8]{inputenc}
\usepackage{amsmath, amssymb}
\usepackage[margin=1in]{geometry}
\usepackage{graphicx}
\newcommand{\pandocbounded}[1]{#1}
% Agregar esto al preámbulo de tu master_template.tex
\newenvironment{answerlist}%
  {\renewcommand{\labelenumii}{(\alph{enumii})}\begin{enumerate}}%
  {\end{enumerate}}

% --- THE MASTER SWITCH TO HIDE SOLUTIONS ---
\usepackage{comment}
\excludecomment{solution} % Tells LaTeX to ignore the solution blocks entirely

% --- Setup for the exams package ---
\newcommand{\exinput}[1]{\input{#1}}
\newenvironment{question}{\item }{} % Formats the questions nicely in your numbered list

\begin{document}

\begin{center}
\large \textbf{An\'alisis Estad\'istico de Datos} \\
\large Examen Argumentativo
\end{center}

Nombre: \underline{\hspace*{3in}} \quad Matr\'icula:\underline{\hspace*{2in}} \\

\noindent \textbf{Instrucciones:} Muestra c\'omo obtienes tu respuesta. Puedes usar Excel 
o Minitab para tus cálculos de probabilidades o 
percentiles, pero no uses otros archivos o hagas uso de la IA. Respeta el 
c\'odigo de \'etica del curso y del Tecnol\'ogico 
de Monterrey.

\begin{enumerate}
    \exinput{exercises}
\end{enumerate}

\end{document}
